%% Chapter 14 ----------------------------------------------------------------
\chapter{Solar-Event Viewers}
\label{ch:solar-events}
\index{solar-event viewers}

The viewers in the \menu{Solar Events} menu place a radio event in its wider context: the
X-ray flare, energetic protons, the geomagnetic response and any associated coronal mass
ejection. Each viewer is an independent window that fetches its data from the archive
listed in Appendix~\ref{app:data-sources}. The viewers can also be started from source as
standalone programs (Section~\ref{sec:from-source}).

\section{Common controls}

The GOES and geomagnetic viewers share a common layout:
\begin{itemize}
  \item a time-range selector with \gui{Start} and \gui{End} rows (date, and hour, minute or
    three-hour slot) and, where applicable, the spacecraft;
  \item a \gui{Preset} list with \gui{Apply} for common ranges;
  \item a plot button that fetches and plots the data, with a progress indicator;
  \item \gui{Save Plot} (PNG) and \gui{Save Data}\index{Burst Analyzer!Excel export} (CSV);
  \item a reset or clear button; and
  \item an information panel summarizing the plotted data or a selected interval.
\end{itemize}

\section{Synchronizing the time window}
\label{sec:sync-time}
\index{time synchronization}

\menu{Solar Events > Sync Current Time Window}\index{time synchronization} sends the time range of the spectrum loaded in
the main window to every open context window: the GOES X-ray and proton viewers, the Dst and
Kp viewers and the STEREO/SWAVES dialog plot the new range; the SOHO/LASCO\index{SOHO/LASCO} CME catalog
searches the day of the middle of the range; and the SunPy Explorer takes the range as its
query window. The status bar reports which windows were updated, or that none were open.

\section{GOES X-ray flux}
\label{sec:goes-xrs}
\index{GOES!X-ray flux}

\menu{Solar Events > Flares > GOES X-Ray Flux} opens the \gui{GOES XRS Plotter}
(Figure~\ref{fig:goes-xrs}), which plots the soft X-ray flux measured by the GOES X-Ray
Sensor in its short (XRS-A, 0.05--0.4~nm) and long (XRS-B, 0.1--0.8~nm) channels.

\screenshot{space-weather/goes_xrs}{The GOES XRS Plotter with a selected flare interval.}{fig:goes-xrs}{The GOES XRS Plotter showing XRS-A and XRS-B on a logarithmic flux axis with flare-class levels, a dragged selection over a flare, and the Flare Information panel filled in.}

\begin{description}
  \item[Start, End, Use] the UTC range and the spacecraft (GOES-8 to GOES-19; default
    GOES-16). \gui{Preset: Whole day (00:00–23:59)} selects the whole of the start day.
  \item[Plot XRS Data] downloads the daily science files from NOAA and plots both channels
    on a logarithmic flux axis.
  \item[Selection] drag across the plot to select a flare interval. The
    \gui{Flare Information}\index{GOES!flare information} panel then shows the \gui{Selection window}, the
    \gui{Peak X-ray Flux (W/m²)} and \gui{Peak Time (UTC)} of the long channel, the
    \gui{Rise Time} (selection start to peak), the \gui{Decay Time} (peak to selection end),
    the \gui{Flare Duration} and the \gui{GOES Class}\index{GOES!flare class} derived from the peak flux.
  \item[Save Plot, Save Data, Reset] save the figure as PNG, save the fluxes as CSV, or clear
    the plot and selection.
\end{description}

\subsection{GOES overlay on the spectrum}
\label{sec:goes-overlay}
\index{GOES!overlay}

\menu{Solar Events > GOES Overlay > Short(XRS-A)}\index{GOES!overlay} and \menu{Long(XRS-B)} draw the GOES
X-ray curves directly on the dynamic spectrum in the main window
(Figure~\ref{fig:goes-overlay}), with a flare-class guide (A, B, C, M, X) on the right. The
overlay does not modify the spectrum data. The application chooses a GOES satellite
appropriate for the date automatically, falling back across legacy and modern satellites
when a file is unavailable. The overlay is included in exported figures.

\screenshot{space-weather/goes_overlay}{GOES X-ray curves overlaid on a dynamic spectrum.}{fig:goes-overlay}{Main window with a dynamic spectrum and both GOES Overlay channels enabled, showing the right-hand flare-class guide (A/B/C/M/X).}

\section{GOES proton flux}
\label{sec:goes-sep}
\index{GOES!proton flux}\index{solar energetic particles}

\menu{Solar Events > Energetic Particles > GOES SEP Proton Flux} opens the
\gui{GOES SEP Proton Flux Plotter} (Figure~\ref{fig:goes-sep}), which plots proton fluxes
from the GOES Solar and Galactic Proton Sensor (SGPS) archive at NOAA over multi-day UTC
ranges.

\screenshot{space-weather/goes_sep}{The GOES SEP Proton Flux Plotter.}{fig:goes-sep}{The GOES SEP Proton Flux Plotter for a multi-day range during an SEP event, with the ~10 MeV and ~100 MeV channels plotted, a selected interval, and the SEP Analysis panel filled in.}

\begin{description}
  \item[Start, End, Use] the UTC range and the spacecraft: \gui{Auto (GOES-19 -> GOES-16)}
    tries the newest spacecraft first, or a specific GOES satellite can be chosen.
  \item[Preset] \gui{Last 24 hours}, \gui{Last 72 hours}, \gui{Last 7 days} or
    \gui{Current UTC day}.
  \item[Plot SEP Proton Flux] fetches and stitches the daily files and plots the channels
    closest to about 10~MeV and 100~MeV, with hover readouts.
  \item[Selection] drag across the plot to measure an event.
  \item[SEP Analysis] the spacecraft, channels, number of samples and source files, the
    peaks of both channels in the plot, and for the selection its window, peak (low
    channel), peak time, rise time, decay time and duration.
  \item[Save Plot, Save Data, Clear Selection] save the figure as PNG, save the stitched
    flux table as CSV, or clear the selection.
\end{description}

\section{Kyoto Dst index}
\label{sec:dst}
\index{Dst index}

\menu{Solar Events > Geomagnetic > Kyoto Dst Index} opens the \gui{Dst Index} viewer
(Figure~\ref{fig:dst}) for the hourly Dst index from the World Data Center for
Geomagnetism, Kyoto \parencite{nose2015}. Dst measures the strength of the ring current;
large negative values indicate geomagnetic storms.

\screenshot{space-weather/dst_index}{The Kyoto Dst Index viewer.}{fig:dst}{The Dst Index viewer for a storm interval, showing the Dst curve with the Quiet, −50, −100 and −200 nT guide lines and the Dst Summary panel.}

\begin{description}
  \item[Start, End] date and hour. \gui{Preset}: \gui{Last 24 hours}, \gui{Last 7 days},
    \gui{Current month} or \gui{Previous month}.
  \item[Plot Dst Data] fetches and plots the index with storm-level guides at 0 (quiet),
    \(-50\), \(-100\) and \(-200\)~nT.
  \item[Dst Summary] the number of samples, the data sources (final, provisional or
    real-time), and the minimum, maximum (each with its time) and mean Dst.
  \item[Save Plot, Save Data, Reset] PNG figure, CSV table, or clear.
\end{description}

\section{GFZ Kp index}
\label{sec:kp}
\index{Kp index}

\menu{Solar Events > Geomagnetic > GFZ Kp Index} opens the \gui{Kp Index} viewer
(Figure~\ref{fig:kp}) for the three-hourly planetary Kp index from GFZ Potsdam
\parencite{matzka2021}.

\screenshot{space-weather/kp_index}{The GFZ Kp Index viewer.}{fig:kp}{The Kp Index viewer for a storm interval, showing colored Kp bars with the Kp 5, 7 and 8 guide lines and the Kp Summary panel.}

\begin{description}
  \item[Start, End] date and three-hour slot (\gui{3h Slot}). \gui{Preset}:
    \gui{Last 24 hours}, \gui{Last 7 days}, \gui{Current month} or \gui{Previous month}.
  \item[Plot Kp Data] plots the index as bars colored by activity (quiet below 5, storm from
    5, severe from 7), with guide lines at Kp~5, 7 and 8.
  \item[Kp Summary] the number of intervals, the maximum, minimum and mean Kp, and the
    number of storm intervals (\(\mathrm{Kp}\geq5\)).
  \item[Save Plot, Save Data, Reset] PNG figure, CSV table, or clear.
\end{description}

\section{SOHO/LASCO CME catalog}
\label{sec:lasco-catalog}
\index{CME catalog}\index{SOHO/LASCO}

\menu{Solar Events > CMEs > SOHO/LASCO CME Catalog} opens the
\gui{SOHO/LASCO CME Catalog Tool} (Figure~\ref{fig:lasco-catalog}), which reads the CDAW\index{CDAW}
catalog of CMEs observed by SOHO/LASCO \parencite{yashiro2004}.

\screenshot{space-weather/lasco_cme_catalog}{The SOHO/LASCO CME Catalog Tool.}{fig:lasco-catalog}{The SOHO/LASCO CME Catalog Tool after Search for a day with several CMEs, with one CME selected, its details shown at the right and the movie status panel with the Play Movie, Open in Browser and Copy URL buttons.}

\begin{description}
  \item[Year, Month, Day, Search] list the CMEs of the chosen day. A progress bar shows the
    retrieval.
  \item[CME table] one row per CME with the columns \gui{Date and Time}, \gui{Central PA},
    \gui{Angular Width}, \gui{Linear Speed}, \gui{Accel}, \gui{Mass},
    \gui{Kinetic Energy}, \gui{MPA} (measurement position angle) and \gui{Remarks}.
    Clicking a row shows its details beside the table.
  \item[Play Movie] plays the LASCO movie associated with the selected CME (double-clicking a
    row does the same) in a viewer with \gui{Back}, \gui{Forward}, \gui{Reload},
    \gui{Open in Browser}, \gui{Open Raw Movie} and \gui{Copy URL}.
  \item[Open in Browser, Copy URL] open the CDAW page of the selected CME, or copy its
    address.
\end{description}
